% Convenciones y comandos compartidos con ../standards/.
\documentclass[letterpaper,12pt]{article}
\usepackage[T1]{fontenc}
\usepackage[utf8]{inputenc}
\usepackage[spanish]{babel}
\spanishdecimal{.}
\usepackage[margin=2.5cm]{geometry}
\usepackage[sc]{mathpazo}
\usepackage{amsmath,amssymb,amsthm}
\usepackage{xcolor}
\usepackage[hidelinks]{hyperref}

% Usar standards en el entorno del autor y la copia incluida al descargar.
\IfFileExists{../standards/definiciones.tex}{%
  \input{../standards/definiciones.tex}%
}{%
  %Si no esta usando el iop package
%\input{definiciones}
%Comandos mios generales
\newcommand{\fig}[1]{fig.\ref{#1}}
\newcommand{\eqn}[1]{eqn.\eqref{#1}}
\newcommand{\tab}[1]{table \ref{#1}}
\newcommand{\medline}{\begin{center}\line(1,0){470}\end{center}}
\newcommand{\suspensivos}{{\color{blue} (.....)}}
\newcommand{\anotacion}[1]{{\color{red} (.. #1 ...)}}
\newcommand{\ie}{\textit{i. e.}}
%Quantum Mechanics Commands
\newcommand{\ket}[1]{{\vert #1 \rangle}}
\newcommand{\bra}[1]{{\langle #1 \vert}}
\newcommand{\proj}[2]{{\vert #1 \rangle \langle #2 \vert}}
\newcommand{\NoM}{{\cal N}}
%Lugares
\newcommand{\unam}{Universidad Nacional Aut\'onoma de M\'exico, M\'exico, D.F., M\'exico}
\newcommand{\icf}{Instituto de Ciencias F\'{\i}sicas, \unam}
\newcommand{\ifunam}{Instituto de F\'{\i}sica, \unam}
\newcommand{\ifunamint}{Instituto de F\'{\i}´ısica, Universidad Nacional Aut\'onoma de M\'exico, M\'exico D.F. 01000, M\'exico}
\newcommand{\fcen}{Departamento de F\'isica ``J. J. Giambiagi'', FCEN, Universidad de Buenos Aires, 1428 Buenos Aires, Argentina}
%%%%%%%%%%% Comandos de Carlos %%%%%%%%%%%%%%%%%%%%%%%%
% enges
%%%%%%%%%%%%%
%%%%%%%%%%%
\newif\ifenglish
\newif\ifespanol
\newif\ifshort
\newif\iflong
\shorttrue
\englishtrue
\newcommand{\enges}[2]{\ifenglish #1 \fi \ifespanol #2 \fi}
%%%%%%
%%%%%
%
}
\englishfalse
\espanoltrue
% Comandos adicionales para estas notas.
\newcommand{\mcE}{\ensuremath{\mathcal{E}}}
\newcommand{\mcL}{\ensuremath{\mathcal{L}}}
\newcommand{\mcT}{\ensuremath{\mathcal{T}}}
\DeclareMathOperator{\tr}{tr}
\DeclareMathOperator{\id}{id}
\newtheorem{teorema}{Teorema}[section]
\renewcommand{\proofname}{Demostración}
\numberwithin{equation}{section}

\title{Notas de mecánica e información cuántica}
\author{David Dávalos}
\date{}

\begin{document}
\maketitle
\begin{center}
\small Versión en desarrollo.\par
\textcopyright\ 2026 David Dávalos.
\href{https://creativecommons.org/licenses/by/4.0/deed.es}{CC BY 4.0}:
se permite compartir y adaptar con atribución e indicación de cambios.
\end{center}

\section*{Notación y convenciones}
Estas convenciones se emplearán a lo largo de las notas. Salvo indicación
contraria, trabajaremos con espacios complejos de dimensión finita.

\paragraph{Espacios y operadores.}
Los espacios de Hilbert se denotan por $\mathsf H$, $\mathsf M$, etc.; los
subíndices identifican sistemas. El espacio de un sistema compuesto es el
producto tensorial de los espacios de sus partes.
Escribimos $\mathcal{B}(\mathsf H)$ para los operadores acotados,
$\mathcal{L}(\mathsf H)$ para los operadores lineales y
$\mathcal{T}(\mathsf H)$ para los operadores de clase traza. En dimensión
finita estos tres espacios coinciden. Los observables se representan por
operadores autoadjuntos $A$, $B$, $G$, etc.; reservamos $E$ para los efectos
de una POVM: $E_a\geq0$ y $\sum_a E_a=\id$.

\paragraph{Estados y notación de Dirac.}
Los vectores de estado se escriben como $\ket{\Psi}$ y sus duales como
$\bra{\Psi}$; en expresiones con muchos estados usaremos minúsculas,
como $\ket{\psi}$. El producto interno $\langle\phi|\psi\rangle$ es
lineal en el ket y antilineal en el bra. Los vectores que representan
estados estarán normalizados. Una fase global no cambia el estado físico:
$\ket{\psi}$ y $e^{i\theta}\ket{\psi}$ representan el mismo rayo.
El operador $\proj{\psi}{\phi}$ es un producto exterior.

Denotamos las matrices de densidad por $\rho$, $\sigma$, $\psi$, etc., y
su conjunto por
\[
  \mathcal{S}(\mathsf H)
  =\{\rho\in\mathcal{T}(\mathsf H):\rho\geq0,\ \tr\rho=1\}.
\]
En particular, $\psi:=\proj{\psi}{\psi}$ denota el estado puro asociado
al ket $\ket{\psi}$. Así, el contexto distingue el operador $\psi$ del
vector $\ket{\psi}$. Abreviamos
$\ket{\psi}\ket{\phi}=\ket{\psi}\otimes\ket{\phi}$ cuando el orden
de los sistemas esté claro.

\paragraph{Promedios y vectores.}
Todo valor esperado se escribe como $\langle A\rangle$; cuando convenga
explicitar el estado, usamos
\[
  \langle A\rangle_\rho=\tr(\rho A),
  \qquad
  \langle A\rangle_\psi=\bra{\psi}A\ket{\psi}.
\]
Los vectores de Bloch y, en general, los vectores de promedios se denotan
por $\vec r$. Para un qubit,
$\rho=\tfrac12(\id+\vec r\cdot\vec\sigma)$, donde
$\vec\sigma=(\sigma_x,\sigma_y,\sigma_z)$ son las matrices de Pauli.

\paragraph{Canales, instrumentos y generadores.}
Los canales cuánticos se denotan por $\mathcal{E}$, $\mathcal{M}$, etc.
Son aplicaciones lineales completamente positivas que preservan la traza;
su acción se escribe con corchetes: $\mathcal{E}[\rho]$.
Reservamos $\mathcal{I}$ para instrumentos cuánticos. Cada operación
$\mathcal{I}_a$ es completamente positiva y no aumenta la traza;
$\tr\mathcal{I}_a[\rho]$ es la probabilidad del resultado $a$.
La suma $\sum_a\mathcal{I}_a$ es un canal y, por tanto, preserva la traza.
Los generadores de Lindblad se escriben al mismo nivel tipográfico:
$\mathcal{L}[\rho]$. Los subíndices pueden indicar parámetros, como
$\mathcal{L}_{H,\beta}$, donde $H$ denota un hamiltoniano y
$\mathsf H$ un espacio de Hilbert.

\paragraph{Estado de Choi.}
Reservamos $\tau$ para matrices de Choi y empleamos por defecto la
versión normalizada, también llamada estado de Jamio\l kowski. Para un
canal $\mathcal{E}$ con espacio de entrada de dimensión $n$, fijamos una
base ortonormal $\{\ket{j}\}_{j=0}^{n-1}$ y definimos
\[
  \ket{\Omega}=\frac{1}{\sqrt n}\sum_{j=0}^{n-1}\ket{j}\ket{j},
  \qquad \omega=\proj{\Omega}{\Omega},
  \qquad \tau_{\mathcal{E}}=(\id_n\otimes\mathcal{E})[\omega].
\]
Aquí $\id_n$ es el canal identidad sobre el sistema de referencia;
$\tr\tau_{\mathcal{E}}=1$. Esta representación depende de la base elegida.
Las referencias bibliográficas se numeran por orden de aparición y se
citan, por ejemplo, como ``Ref.~\cite{preskill}''.

\section{Teorema de no clonación}
Clonar un estado significa producir dos copias exactas, conservando el
original. El segundo sistema comienza en un estado fijo $\ket{0}$,
que sirve como soporte de la copia. Incluimos además un espacio
$\mathsf M$ con todos los grados de libertad auxiliares de la máquina.
La formulación con estos grados de libertad puede consultarse en la
Ref.~\cite{preskill}, sección 4.2.3.

\begin{teorema}[No clonación con grados de libertad de la máquina]
Sean $\ket{\psi},\ket{\phi}\in\mathsf H$ dos vectores normalizados
que representan estados puros distintos y no ortogonales; es decir,
\[
  0<|\langle\psi|\phi\rangle|<1.
\]
No existen una misma unitaria $U$ sobre
$\mathsf H\otimes\mathsf H\otimes\mathsf M$ y estados iniciales
fijos y normalizados $\ket{0}\in\mathsf H$, $\ket{e}\in\mathsf M$
que satisfagan simultáneamente
\begin{align}
  U\bigl(\ket{\psi}\ket{0}\ket{e}\bigr)
    &=\ket{\psi}\ket{\psi}\ket{e'_\psi},\label{eq:clona-psi}\\
  U\bigl(\ket{\phi}\ket{0}\ket{e}\bigr)
    &=\ket{\phi}\ket{\phi}\ket{e'_\phi},\label{eq:clona-phi}
\end{align}
aun permitiendo que los estados finales normalizados de la máquina,
$\ket{e'_\psi}$ y $\ket{e'_\phi}$, sean diferentes.
En particular, si $\dim\mathsf H\geq2$, no existe una máquina que clone
exactamente todos los estados puros.
\end{teorema}

\begin{samepage}
\begin{proof}
Una unitaria conserva los productos internos. Al tomar el producto
interno entre las entradas y entre las salidas de
\eqref{eq:clona-psi}--\eqref{eq:clona-phi}, obtenemos
\begin{equation}
  \langle\psi|\phi\rangle
    =\langle\psi|\phi\rangle^2
      \langle e'_\psi|e'_\phi\rangle.
\end{equation}
Sea $s=|\langle\psi|\phi\rangle|$. Por la desigualdad de
Cauchy--Schwarz y la normalización de los estados de la máquina,
\[
  s=s^2|\langle e'_\psi|e'_\phi\rangle|\leq s^2.
\]
Esto es imposible para $0<s<1$, pues entonces $s^2<s$.
Finalmente, todo espacio de dimensión al menos dos contiene un par de
estados de este tipo: si $\ket{0},\ket{1}$ son ortonormales, basta
elegir $\ket{0}$ y $(\ket{0}+\ket{1})/\sqrt2$.
\end{proof}
\end{samepage}

\paragraph{Qué estados sí se pueden clonar.}
La restricción anterior permite únicamente $s=0$ (estados ortogonales)
o $s=1$ (el mismo estado físico). Recíprocamente, cualquier familia de
estados puros mutuamente ortogonales $\{\ket{j}\}$ puede clonarse:
la asignación
\[
  \ket{j}\ket{0}\ket{e}\longmapsto\ket{j}\ket{j}\ket{e'}
\]
con $\ket{e'}$ fijo y normalizado, lleva una familia ortonormal a otra.
Completando ambas a bases
ortonormales del espacio total, se extiende a una unitaria.
Por tanto, una familia de estados puros distintos admite clonación
exacta y determinista si y solo si sus estados son mutuamente ortogonales.

\paragraph{Por qué incluir la máquina basta.}
El estado inicial de la máquina es fijo: no contiene información previa
acerca de cuál estado se recibe. Su estado final puede depender libremente
de la entrada. Además, si las dos copias son puras y exactas, el estado
conjunto de las copias es puro y no puede estar entrelazado con la máquina;
de ahí la forma factorizada de las salidas.
Una operación cuántica determinista general puede realizarse mediante una
unitaria sobre un sistema ampliado y el descarte de auxiliares. Si la
máquina comienza en un estado mixto, se incluye también un sistema que lo
purifique. Por ello, permitir auxiliares, mediciones sin postselección o
descartar la máquina no evita el teorema. La afirmación se refiere a
clonación \emph{exacta y determinista}; no excluye copias aproximadas ni
protocolos probabilísticos para ciertas familias de estados.

\bibliographystyle{unsrt}
\IfFileExists{../standards/labibliografia.bib}{%
  \bibliography{../standards/labibliografia}%
}{%
  \bibliography{support/labibliografia}%
}

\end{document}
